% Appendix A

\chapter{Additional Methodological Details}

\label{AppendixA}


\section{Baseline Elasticity Simulation}

For each holdout reference-week, the predictive LightGBM baseline is evaluated over the operational discount grid from 20\% to 70\%. All contextual variables are held fixed, while the discount variable (\texttt{discount}) is changed. This creates a simulated demand curve for the same reference-week under alternative markdown decisions.

The predicted weekly sales at each discount are normalized by the prediction at the 20\% reference discount:

\begin{equation}
\text{sales multiplier}(d) =
\frac{\widehat{\text{sales}}(d)}
{\widehat{\text{sales}}(20\%)}.
\end{equation}

The break-even multiplier is the increase in units sold required to compensate for the lower discounted price:

\begin{equation}
\text{break-even}(d) =
\frac{P(20\%)}
{P(d)}
=
\frac{1-20\%}{1-d},
\end{equation}

where $P(d)$ is discounted price at discount $d$.

A simulated discount is revenue-improving relative to 20\% only when the predicted sales multiplier is above the break-even multiplier. This diagnostic is used to identify whether the predictive baseline produces economically reasonable curves or instead generates boundary-oriented recommendations.

\section{Hyperparameters}

The LightGBM nuisance models used in the DML step are trained to estimate the conditional expectation of the outcome and the treatment given the pre-pricing covariates. The same configuration is used for both nuisance models. Cross-fitting is performed using leave-one-season-out folds to reduce temporal leakage.

\begin{table}[h]
\centering
\small
\begin{tabular}{ll}
\toprule
\textbf{Hyperparameter} & \textbf{Value} \\
\midrule
Objective & regression \\
Maximum boosting iterations & 1000 \\
Learning rate & 0.01 \\
Number of leaves & 127 \\
Minimum child samples & 8 \\
Subsample fraction & 0.8 \\
Column subsample by tree & 0.8 \\
L2 regularization & 0.03 \\
Early stopping rounds & 20 \\
Random seed & 42 \\
Cross-fitting & leave-one-season-out \\
\bottomrule
\end{tabular}
\caption{Hyperparameters of the LightGBM nuisance models used in the DML residualization step.}
\label{tab:dml-hyperparams}
\end{table}

No monotonic constraint is imposed at the nuisance stage: the nuisance models estimate conditional means, while the causal price effect is captured in the final stage. Leave-one-season-out cross-fitting blocks temporal leakage and avoids predicting residuals on observations used to train the corresponding nuisance model.

The Causal Forest is used as the heterogeneous final stage after residualization. It receives the residualized outcome, the residualized treatment, and the pre-pricing covariates used as heterogeneity features.

\begin{table}[h]
\centering
\small
\begin{tabular}{ll}
\toprule
\textbf{Hyperparameter} & \textbf{Value} \\
\midrule
Number of trees & 1000 \\
Minimum samples per leaf & 80 \\
Random seed & 42 \\
Parallel jobs & -1 \\
Honest splitting & enabled \\
Inference enabled & yes \\
\bottomrule
\end{tabular}
\caption{Hyperparameters of the Causal Forest final stage.}
\label{tab:cf-hyperparams}
\end{table}

The minimum leaf size is set conservatively to reduce noisy reference-level treatment effects. This is important because the final output is not only used for average interpretation, but also for individual reference-week elasticity curves. The nuisance models are not internal to the forest; they are fitted separately with LightGBM and their residuals are passed to \texttt{CausalForest}.
